\ifdefined\gkver\else\def\gkver{student}\fi
\documentclass[\gkver]{gaokaozhenti}
\begin{document}

\chapter{2002年上海春季理综}
%% paperType: 综合能力测试物理部分

\begin{enumerate}
\item
%% number: 13
%% paperName: 2002年上海春季理综第 13 题 3 分
%% typeId: 1
%% type: 单选题
%% chapter: 运动和力
%% point: 牛顿第一定律
%% method: 无
%% score: 3
%% degree: 650
%% duplicateId: 0
%% body:
根据牛顿运动定律，以下选项中正确的是\xzanswer{C}

\fourchoices[answer=C]
{人只有在静止的车厢内，竖直向上高高跳起后，才会落在车厢的原来位置}
{人在沿直线匀速前进的车厢内，竖直向上高高跳起后，将落在经起点的后方}
{人在沿直线加速前进的车厢内，竖直向上高高跳起后，将落在起跳点的后方}
{人在沿直线减速前进的车厢内，竖直向上高高跳起后，将落在起跳点的后方}

%% answer: C
%% memo:
\memoanswer{
本题原卷及材料中未提供详解，待补充。
}

\item
%% number: 14
%% paperName: 2002年上海春季理综第 14 题 3 分
%% typeId: 1
%% type: 单选题
%% chapter: 光学
%% point: 光的干涉
%% method: 无
%% score: 3
%% degree: 650
%% duplicateId: 0
%% body:
牛顿为了说明光的性质，提出了光的微粒说。如今，人们对光的性质已有了进一步的认识。下列四个示意图所表示的实验，能说明光性质的是\xzanswer{B}

\fourpicture[labelprefix={图}, numA=①, numB=②, numC=③, numD=④, h=2.0cm]
{figs/14a.png}
{figs/14b.png}
{figs/14c.png}
{figs/14d.png}

\fourchoices[answer=B]
{①②}
{②③}
{③④}
{②④}

%% answer: B
%% memo:
\memoanswer{
本题原卷及材料中未提供详解，待补充。
}

\item
%% number: 15
%% paperName: 2002年上海春季理综第 15 题 3 分
%% typeId: 1
%% type: 单选题
%% chapter: 磁场
%% point: 电流的磁场
%% method: 无
%% score: 3
%% degree: 850
%% duplicateId: 0
%% body:
右图是一种利用电磁原理制作的充气泵的结构示意图。其工作原理类似打点计时器。当电流从电磁铁的接线柱 $a$ 流入，吸引小磁铁向下运动时，以下选项中正确的是\xzanswer{D}

\onepicture[h=2.6cm]{figs/15.jpg}

\fourchoices[answer=D]
{电磁铁的上端为 N 极，小磁铁的下端为 N 极}
{电磁铁的上端为 S 极，小磁铁的下端为 S 极}
{电磁铁的上端为 N 极，小磁铁的下端为 S 极}
{电磁铁的上端为 S 极，小磁铁的下端为 N 极}

%% answer: D
%% memo:
\memoanswer{
本题原卷及材料中未提供详解，待补充。
}

\item
%% number: 16
%% paperName: 2002年上海春季理综第 16 题 3 分
%% typeId: 1
%% type: 单选题
%% chapter: 气体性质
%% point: 玻意耳定律
%% method: 无
%% score: 3
%% degree: 750
%% duplicateId: 0
%% body:
右图为充气泵气室的工作原理图。设大气压强为 $p_{0}$，气室中的气体压强为 $p$，气室通过阀门 $\mathrm{K}_{1}$、$\mathrm{K}_{2}$ 与空气导管相连接。以下选项中正确的是\xzanswer{C}

\onepicture[h=3.0cm]{figs/16.jpg}

\fourchoices[answer=C]
{当橡皮碗被拉伸时，$p>p_{0}$，$\mathrm{K}_{1}$ 关闭，$\mathrm{K}_{2}$ 开通}
{当橡皮碗被拉伸时，$p<p_{0}$，$\mathrm{K}_{1}$ 关闭，$\mathrm{K}_{2}$ 开通}
{当橡皮碗被压缩时，$p>p_{0}$，$\mathrm{K}_{1}$ 关闭，$\mathrm{K}_{2}$ 开通}
{当橡皮碗被压缩时，$p<p_{0}$，$\mathrm{K}_{1}$ 关闭，$\mathrm{K}_{2}$ 开通}

%% answer: C
%% memo:
\memoanswer{
本题原卷及材料中未提供详解，待补充。
}

\end{enumerate}

\end{document}
